\documentclass[11pt]{article}

\usepackage[utf8]{inputenc}
\usepackage[T1]{fontenc}
\usepackage{lmodern}
\usepackage[margin=1in]{geometry}
\usepackage{microtype}
\usepackage{parskip}
\usepackage[hidelinks]{hyperref}

\title{\vspace{-2em}\textbf{Teaching Statement}\\[0.3em]\large Alexandre Pavlov}
\author{}
\date{}

\begin{document}
\maketitle
\vspace{-3em}

\begin{center}
\small Université de Montréal \hspace{0.5em}\textbar\hspace{0.5em}
\href{mailto:alexandre.pavlov@umontreal.ca}{alexandre.pavlov@umontreal.ca} \hspace{0.5em}\textbar\hspace{0.5em}
\href{https://alexandrepavlov.com/}{alexandrepavlov.com}
\end{center}
\vspace{0.5em}

\textbf{Teaching philosophy.}
Two convictions guide my teaching. First, economics is learned by reasoning, not memorization: a student who can rebuild the logic behind the aggregate demand curve will outperform one who has memorized its shape. Second, reasoning is built through dialogue rather than lecture. Students come to understand a concept when they state it in their own words and have their logic corrected, refined, or extended. These convictions shape how I structure my classes and how I respond to students inside and outside them.

\textbf{Experience.}
I have taught economics for seven years, at McGill University and the Université de Montréal, from first-year lectures to the PhD core. Since 2022 I have been the instructor of record for three first-year courses at the Université de Montréal (macroeconomics, economics and globalization, and introduction to economics), teaching sections of fifty to one hundred students. As a teaching assistant I have supported courses from undergraduate econometrics and the third-year applied research workshop to the master's course in macroeconometrics and the first-year PhD course in macroeconomics. Students rate my teaching 3.79 out of 4.00 on average, against a departmental average of about 3.4.

\textbf{Building intuition before formalism.}
Students arrive in macroeconomics expecting to memorize a sequence of curves and equations. I structure lectures so that the formalism is the destination, never the starting point. I open each topic with a current policy question: how to design a carbon price, what central-bank tightening will do to output and employment, whether fiscal stimulus is worth its cost in a recession. The question gives the model a job to do. Students see why we need the model before we write it down, and they retain the mechanics better as a result. Opening this way also lets me draw on my research in climate macroeconomics and public finance, so first-year students meet debates the discipline has not yet settled.

\textbf{Dialogue and the productive use of wrong answers.}
When a student offers an incorrect answer, I do not correct it and move on: I ask them to walk through their reasoning, and we follow it together until we find the step where the logic broke. The point is to make thinking visible, to the student and to the rest of the class. Over a semester this changes the room. Students who said nothing in the first weeks volunteer answers by midterm, because a wrong step is treated as part of learning rather than as failure. Treating error this way also surfaces misunderstandings that a one-way lecture would hide until the exam.

\textbf{Adapting to diverse backgrounds.}
A first-year class at the Université de Montréal mixes students with very different preparation. Some have completed advanced calculus; others have not seen a derivative in years. Some are economics majors; others arrive from political science, history, or environmental studies. Many follow the course in their second or third language. I design lectures and problem sets so that the core material is within reach of every student in the room, and I leave paths open for those who want more: optional readings, harder end-of-set questions, longer conversations in office hours. Having taught in French at the Université de Montréal and in English at McGill, I know where technical vocabulary fails to travel across languages, and I write key terms on the board so that students working in their second language do not lose the thread.

\textbf{Feedback as an ongoing input.}
I treat student feedback as an input during the semester, not only after it. I check in with students informally, after class and in office hours, about what is working and what is not, and I read problem sets and midterms as evidence on my teaching as much as on their learning: when half the class misses the same question, the next lecture revisits that concept from another angle. Pacing, problem-set design, and the balance between intuition and derivation have all moved in response, and my evaluations have improved year over year as a result.

\textbf{Looking ahead.}
I want to build courses that connect my research to the classroom. Few undergraduates ever encounter climate economics formally. I would welcome the chance to design a course on the social cost of carbon, optimal climate policy, and the macroeconomics of the energy transition, taught from the canonical papers and from debates still open in the field. At the graduate level I am prepared to teach core macroeconomics, public finance, and applied climate economics, including the computational tools (Julia, Python, MATLAB, Fortran) that increasingly define applied work in the field. I see teaching and research as continuous: the questions that drive my research (how policy should respond to climate risk, how heterogeneity shapes optimal taxation, how rules and discretion trade off over long horizons) are the questions I want my students to be able to reason about by the end of a semester with me.

\end{document}
